This work tackles an understudied area of the document analysis field: \gls{ZSL}-\gls{DIC}. One of the biggest obstacles of this area is the lack of a specialized dataset. Therefore this work starts by introducing a novel document image dataset, \gls{LA-CDIP}, categorized mainly by the layout information of the document image, which allows for knowledge generalization on a zero-shot setting. This dataset is constructed with a labeling framework assisted by \gls{DL} models that were trained on the dataset itself. This labeling framework works in iterations, where the current loop should be easier to label manually than the previous iteration, allowing the manual labeling to become increasingly easier, as the iterations go by. The resulting dataset is specialized in \gls{ZSL}-\gls{DIC}, providing enough class diversity to allow \gls{DL} models trained from scratch to achieve a high generalization potential.

To use this dataset in the \gls{ZSL}-\gls{DIC}, this work also presents \gls{VDM}, which is a framework that frames the classification problem as a matching problem, using embedding similarity to achieve the \gls{ZSL} paradigm. Therefore, by having a set of ground-truth references, this method can classify any document class in a zero-shot scenario, if correctly trained. To further encourage research with this framework, the dataset is extensively benchmarked by training siamese networks with a wide variety of neural backbones, as well as an ample set of hyperparameters.

In conclusion, this work developed an understudied branch in the \gls{DIC} research, providing a novel dataset and a method, resulting in a classification solution with real-world applications.

\section{Future Works}

There are some known limitations of the work and plans to tackle them. First, as mentioned, the architecture complexity is currently an obstacle, as the dataset is still relatively small. To solve this problem, it is already an ongoing work to keep iterating with the techniques presented here, increasing both the number of classes and number of documents. To further increase the quality of the dataset, \gls{DL}-assisted dataset verification methods are being developed, where the classes are matched against each other to find possible duplicate patterns. This technique will be able to pinpoint possible human mistakes and increase the parallelization of human effort in creating this dataset. Another possible dataset enhancement is to increase the number of document sources by gathering documents from sources other than the RVL-CDIP dataset, increasing the generalization of the trained models.

With regard to the training framework, a possible step to enhance the results is to explore data augmentation techniques, which should alleviate the burden in the dataset labeling process. In the metric learning topic, more advanced techniques, such as tuple mining \cite{schroff_facenet_2015}, are mapped as future experiments, together with exploring the combination of different loss functions and distance functions.